\section{Experimental Design}

\paragraph{Workflow scenario.}
The main benchmark models an enterprise vendor-update workflow. Internal agents coordinate over private memory, shared raw memory, safe-view memory, workspace documents, and inter-agent messages to prepare a vendor-facing project update. The task permits coarse disclosure, such as acknowledging that a budget constraint or internal timing issue exists, but forbids exact budget values, internal incident rationales, private customer/project references, and credential-like markers. The final writer must preserve non-sensitive task utility while withholding raw protected values.

\paragraph{Agents, topologies, and research questions.}
The workflow uses four benchmark roles: a privileged planner, a privileged finance agent, an internal document writer, and an external vendor endpoint. We run the same task, policies, attacks, defenses, and seeds under three runtime topologies. In \emph{chain\_4}, information flows planner $\rightarrow$ finance $\rightarrow$ document writer $\rightarrow$ vendor. In \emph{star\_4}, the planner is a hub connecting the other roles. In \emph{blackboard\_4}, all agents can interact through an all-to-all graph and a shared workspace artifact, increasing fanout. The evaluation asks: RQ1, do final-output and surface defenses miss intermediate propagation? RQ2, does topology affect propagation? RQ3, does FlowFence reduce raw and external leakage while preserving task success? RQ4, can containment be achieved without evaluator-only attack labels? Detailed role, topology, and evidence-matrix tables are in Appendix~\ref{app:details}.

\paragraph{Attacks and defenses.}
The LLM final-writer matrix uses seven attack settings: none, summary-poisoning direct/indirect, workspace-poisoning direct/indirect, and communication-hijack direct/indirect. Direct attacks explicitly request sensitive details; indirect attacks embed the request in operational text. Reserved paraphrase variants are used only in the label-free validation. We compare no defense, static access control (static ACL), prompt filtering, and FlowFence. Static ACL blocks hard forbidden channels and unauthorized raw-recipient cases but does not use semantic, provenance, topology, or fanout reasoning. Prompt filtering uses direct phrase/pattern matching and does not rewrite safe views or track propagation. FlowFence applies policy, provenance, topology, safe-view, quarantine, and propagation-right mediation.

\paragraph{Run matrices and provider configuration.}
The main final-writer experiment uses MiniMax to generate vendor-facing updates from mediated runtime context; MiniMax is not used as the guard, risk scorer, metric evaluator, browser agent, or desktop agent. The main matrix contains $3$ topologies $\times$ $7$ attack settings $\times$ $4$ defenses $\times$ $3$ seeds. The label-free validation contains a deterministic matrix with $3$ topologies $\times$ $10$ attacks $\times$ $6$ defenses $\times$ $3$ seeds, plus a targeted MiniMax final-writer matrix with $2$ topologies $\times$ $3$ paraphrase attacks $\times$ $4$ defenses $\times$ $3$ seeds. The committed configuration records temperature $0.0$, max tokens $256$, and timeout $60$ seconds; provider details are reported in Appendix~\ref{app:expdetails}.

\paragraph{Metrics.}
Task success is a deterministic final-output rule/template check for required non-sensitive content, reported separately from privacy leakage. Unauthorized raw leakage counts exact raw-value exposure in runtime event text outside quarantine or blocked events. External leakage counts the subset reaching an external recipient, final output, vendor-send tool, or external message. Cascade size, cascade depth, and privilege reach are computed over the full event graph. Evaluator-label use counts defense decisions whose metadata indicates use of evaluator-only attack labels. These metrics measure runtime propagation rather than only final-answer safety.
